\section*{第 \ref{ch:b1:gene-expression} 章——基因表达：转录与翻译}

\begin{solution}{exo:b1:gene-expression:1}
\omterm{def:b1:nucleic-acids:chain}{RNA} 为 $5'$-AUGCCUAAG-$3'$；编码链为 $5'$-ATGCCTAAG-$3'$。
\end{solution}

\begin{solution}{exo:b1:gene-expression:2}
加上 $5'$ 帽（修饰的鸟\omterm{def:b1:nucleic-acids:nucleotide}{嘌呤}）、$3'$ 端 poly-A 尾，以及
通过剪接去除\omterm{def:b1:genomes:gene}{内含子}。
\end{solution}

\begin{solution}{exo:b1:gene-expression:3}
AUG CCG AAA GUU UGA：Met-Pro-Lys-Val，随后遇终止\omterm{def:b1:gene-expression:code}{密码子}。
\end{solution}

\begin{solution}{exo:b1:gene-expression:4}
A 位：进入的带电 tRNA 被接纳并检验。P 位：携带
延伸中肽链的 tRNA；\omterm{def:b1:proteins:peptide}{肽键}在其肽链与
A 位\omterm{def:b1:proteins:aminoacid}{氨基酸}之间形成。E 位：正在离去的空载 tRNA。
\end{solution}

\begin{solution}{exo:b1:gene-expression:5}
编码区 $2400 - 600 = 1800$ 个\omterm{def:b1:nucleic-acids:nucleotide}{核苷酸}：599 个残基（600 个\omterm{def:b1:gene-expression:code}{密码子}，
其中一个为终止\omterm{def:b1:gene-expression:code}{密码子}）。所需时间 $599/4 = \qty{150}{s}$。同时工作的\omterm{def:b1:gene-expression:ribosome}{核糖体} $1800/90 = 20$ 个。
\end{solution}

\begin{solution}{exo:b1:gene-expression:6}
信使 \omterm{def:b1:nucleic-acids:chain}{RNA} 从固定的起始位点起每三个碱基读取一次；多出一个或两个
额外碱基会使下游每个\omterm{def:b1:gene-expression:code}{密码子}错位，其后的\omterm{def:b1:proteins:peptide}{蛋白质}
全是无意义序列，通常在提前出现的终止处结束。三个额外碱基
增加一个\omterm{def:b1:gene-expression:code}{密码子}并恢复\omterm{def:b1:gene-expression:code}{阅读框}：\omterm{def:b1:proteins:peptide}{蛋白质}多出一个残基，
两次插入之间有少数错误残基，往往仍能发挥功能。
\end{solution}

\begin{solution}{exo:b1:gene-expression:7}
CAG（Gln）变为 UAG（终止）：肽链在第 150 位残基处结束，成为截短的、
几乎肯定无活性的\omterm{def:b1:proteins:peptide}{蛋白质}（无义突变）。CAG 变为 CAA：
仍为 Gln，属沉默突变。CAG 变为 CGG：Gln 变为 Arg，是一次替换
（错义突变），其效应取决于所在位置。
\end{solution}

\begin{solution}{exo:b1:gene-expression:8}
\omterm{def:b1:gene-expression:ribosome}{核糖体}只检验\omterm{def:b1:gene-expression:code}{密码子}与\omterm{def:b1:gene-expression:trna}{反密码子}的配对，它看不到
\omterm{def:b1:proteins:aminoacid}{氨基酸}本身。若其 tRNA 上的半胱氨酸被化学转化为
丙氨酸，则\omterm{def:b1:gene-expression:ribosome}{核糖体}会在信使 \omterm{def:b1:nucleic-acids:chain}{RNA} 指定半胱氨酸的每处插入丙氨酸，
因为读取的是 tRNA，而非\omterm{def:b1:proteins:aminoacid}{氨基酸}。因此，氨酰 tRNA 合成酶
负责把每种\omterm{def:b1:proteins:aminoacid}{氨基酸}与正确的 tRNA 配对，才是真正的
翻译者。
\end{solution}

\begin{solution}{exo:b1:gene-expression:9}
$400\times 4 = 1600$ 个键；\omterm{def:b1:gene-expression:ribosome}{核糖体}每个残基的两个 GTP 占其中
一半，合成酶的两个占另一半。
\end{solution}

\begin{solution}{exo:b1:gene-expression:10}
在真核生物中，转录本在\omterm{def:b1:eukaryotic-cell:nucleus}{细胞核}中产生，而\omterm{def:b1:gene-expression:ribosome}{核糖体}
位于\omterm{def:b1:eukaryotic-cell:organelle}{细胞溶胶}中，二者由\omterm{def:b1:eukaryotic-cell:nucleus}{核膜}隔开；转录本在剪接和加帽
完成之前，并不是信使 \omterm{def:b1:nucleic-acids:chain}{RNA}。
这种分隔使加工过程本身成为可能——可变
剪接、输出的调控，以及在阅读前检查信息
是否完整。
\end{solution}

\begin{solution}{exo:b1:gene-expression:11}
四种信使 \omterm{def:b1:nucleic-acids:chain}{RNA}：同时含两个\omterm{def:b1:genomes:gene}{外显子}（增加 190 个\omterm{def:b1:nucleic-acids:nucleotide}{核苷酸}）、仅含\omterm{def:b1:genomes:gene}{外显子} 3
（90）、仅含\omterm{def:b1:genomes:gene}{外显子} 4（100），两个\omterm{def:b1:genomes:gene}{外显子}都不含（0）。\omterm{def:b1:gene-expression:code}{阅读框}是否保持，取决于
增加的长度是否为 3 的倍数：0 和 90 保持\omterm{def:b1:gene-expression:code}{阅读框}；100 和 190 使\omterm{def:b1:gene-expression:code}{阅读框}移位，
因此仅含\omterm{def:b1:genomes:gene}{外显子} 4 或两个\omterm{def:b1:genomes:gene}{外显子}都含的信使 \omterm{def:b1:nucleic-acids:chain}{RNA} 会在\omterm{def:b1:genomes:gene}{外显子} 5 处错读\omterm{def:b1:gene-expression:code}{阅读框}，
并截短\omterm{def:b1:proteins:peptide}{蛋白质}。可变剪接必须遵守\omterm{def:b1:gene-expression:code}{阅读框}，
因此\omterm{def:b1:genomes:gene}{外显子}长度往往是 3 的倍数，
原因正在于此。
\end{solution}

\begin{solution}{exo:b1:gene-expression:12}
任意的：化学上没有任何依据说明 UUU 必须代表 Phe；
这种对应是由合成酶维持的约定。优化的：
第三位是简并性最强的位置，因此落在该位置的错误与突变
往往是沉默的；只相差一个碱基的\omterm{def:b1:gene-expression:code}{密码子}
倾向于编码相似的\omterm{def:b1:proteins:aminoacid}{氨基酸}，因此替换往往较温和
——\omterm{def:b1:gene-expression:code}{遗传密码}使错误造成的危害最小化。固定的：一种合成酶
改变了其专一性，便会同时改变\omterm{def:b1:cell-unit-of-life:cell}{细胞}的每一种\omterm{def:b1:proteins:peptide}{蛋白质}，
即数以千计的\omterm{def:b1:proteins:peptide}{蛋白质}；几乎没有\omterm{def:b1:cell-unit-of-life:cell}{细胞}能承受这种变化，
因此\omterm{def:b1:gene-expression:code}{遗传密码}无法漂移，自所有生物的
共同祖先以来便已固定。
\end{solution}

\begin{solution}{pb:b1:gene-expression:1}
\textbf{1.} $30\,000/30 = \qty{1000}{s}$，约 \qty{17}{min}。
\textbf{2.} $28\,000/30\,000 = \qty{93}{\%}$.
\textbf{3.} $30\,000/2000 = 15$.
\textbf{4.} $2\times 30\,000 = \num{60000}$ 个 \omterm{def:b1:water-small-molecules:atp}{ATP}。
\textbf{5.} $1000/10 = 100$ 个 \omterm{def:b1:nucleic-acids:chain}{RNA} 聚合酶位于该\omterm{def:b1:genomes:gene}{基因}上；
每小时产生 360 个信使 \omterm{def:b1:nucleic-acids:chain}{RNA}。
\textbf{6.} 七个\omterm{def:b1:genomes:gene}{内含子}，平均长 $28\,000/7 = \qty{4000}{}$ 个\omterm{def:b1:nucleic-acids:nucleotide}{核苷酸}。
\textbf{7.} 转录（\qty{17}{min}）；\omterm{def:b1:genomes:gene}{内含子}在
合成的同时即被剪接，最后一个仅再增加一分钟。
\textbf{8.} \omterm{def:b1:genomes:gene}{基因} $30\,000\times 0.34\,\mathrm{nm} = \qty{10}{\micro m}$；信使 \omterm{def:b1:nucleic-acids:chain}{RNA} 长 \qty{0.68}{\micro m}。
\textbf{9.} $500/4 = \qty{125}{s}$.
\textbf{10.} $2000/80 = 25$.
\textbf{11.} 每隔 $80/4 = \qty{20}{s}$ 完成一条肽链，每
小时 180 条。
\textbf{12.} 满负荷生产 $2/\ln 2 = \qty{2.9}{h}$：约 520
个\omterm{def:b1:proteins:peptide}{蛋白质}。
\textbf{13.} $4\times 500 = 2000$ 个 \omterm{def:b1:water-small-molecules:atp}{ATP}。
\textbf{14.} 每个\omterm{def:b1:proteins:peptide}{蛋白质} $60\,000/520 = 115$ 个 \omterm{def:b1:water-small-molecules:atp}{ATP}。
\textbf{15.} 约 2100 个 \omterm{def:b1:water-small-molecules:atp}{ATP}，其中 \qty{95}{\%} 用于翻译。
\textbf{16.} $1 - (1 - 10^{-5})^{1503} \approx 1503\times 10^{-5} =
\qty{1.5}{\%}$.
\textbf{17.} $1 - (1 - 10^{-4})^{500} \approx \qty{5}{\%}$.
\textbf{18.} 转录：$1.5\%\times 0.75\times 0.67 \approx
0.75\%$ 的信使 \omterm{def:b1:nucleic-acids:chain}{RNA} 有缺陷，因而\omterm{def:b1:proteins:peptide}{蛋白质}亦然；翻译：
$5\%\times 0.67 = 3.3\%$；约 \qty{4}{\%} 的分子有缺陷。
\textbf{19.} \omterm{def:b1:proteins:peptide}{蛋白质}错误仅限于一个分子，该分子
很快被降解；信使 \omterm{def:b1:nucleic-acids:chain}{RNA} 的错误会被复制到数百个\omterm{def:b1:proteins:peptide}{蛋白质}中；
复制错误则被每个后代永远继承。
错误的代价，以及值得为之付出的精确度，随着
每往回推一步而升高。
\textbf{20.} 存在 $3000/180 = 17$ 个信使 \omterm{def:b1:nucleic-acids:chain}{RNA}。
\textbf{21.} $17\times 25 = 420$ 个\omterm{def:b1:gene-expression:ribosome}{核糖体}。
\textbf{22.} $\num{3e9}\times 4 = \num{1.2e10}$ 个 \omterm{def:b1:water-small-molecules:atp}{ATP}，每天
$= \qty{2e-14}{mol}$，\qty{1e-9}{J}，即 \qty{1.2e-14}{W}
——每立方微米 \qty{6}{fW}。
\textbf{23.} 每秒 $10^9$ 个 \omterm{def:b1:water-small-molecules:atp}{ATP}，即每天 $\num{8.6e13}$ 个：\omterm{def:b1:proteins:peptide}{蛋白质}
合成在这种更新缓慢的\omterm{def:b1:cell-unit-of-life:cell}{细胞}中仅占其百分之一（在
分裂的细菌中则占一半）。
\textbf{24.} 无成熟的信使 \omterm{def:b1:nucleic-acids:chain}{RNA}：初级转录本在
\omterm{def:b1:eukaryotic-cell:nucleus}{细胞核}中积累，而\omterm{def:b1:proteins:peptide}{蛋白质}则随其信使 \omterm{def:b1:nucleic-acids:chain}{RNA} 降解，
数小时内消失。细菌\omterm{def:b1:proteins:peptide}{蛋白质}，其\omterm{def:b1:genomes:gene}{基因}无\omterm{def:b1:genomes:gene}{内含子}，
不受影响。
\textbf{25.} 每个分子约 2100 个 \omterm{def:b1:water-small-molecules:atp}{ATP}——2000 个用于翻译，
100 个用于其转录本所分担的部分；首个\omterm{def:b1:proteins:peptide}{蛋白质}在
\qty{1000}{s} 的转录后、经一分钟加工与输出，
并加上 \qty{125}{s} 的翻译：约二十分钟。
\end{solution}
